\section{データ}
\label{sec:data}

\subsection{再解析とは何か}

天気の観測は、あるところでは密で、あるところではまばらである。
陸上には気象台もラジオゾンデもあるが、
太平洋の真ん中の高度 \SI{25}{\kilo\metre} を毎日測っている人はいない。
だから「過去のある日の、ある地点の、ある高度の風」を観測だけから並べることはできない。

\textbf{再解析}（reanalysis）はこの穴を埋める方法である。
数値予報モデル---大気の運動方程式を格子上で解くプログラム---を走らせながら、
その時刻に得られた観測（地上観測、ラジオゾンデ、航空機、船舶、衛星）を
次々に取り込んでモデルの状態を修正していく。
これを過去の全期間にわたって、\textbf{同一のモデルと同一の手法で}やり直したものが再解析である。

\begin{whybox}{なぜ「同一の手法で」が重要なのか}
現業の天気予報モデルは毎年改良される。
そのため、現業の解析結果をそのまま何十年も並べると、
気候が変わったのかモデルが変わったのか区別できない。
再解析は全期間を一つのモデルでやり直すので、
少なくともモデル由来の不連続は持ち込まれない。

ただし観測網そのものは時代とともに変わっている（衛星観測は 1979 年以降に本格化した）。
これは再解析でも消せない。本報告が 2016 年以降に限った理由の一つである。
\end{whybox}

したがって再解析の値は\textbf{観測そのものではない}。
観測に強く拘束された、モデルによる最良推定である。
観測が密なところ（日本の陸上、対流圏）では観測に近く、
まばらなところ（高層、外洋）ではモデルの寄与が大きい。

\subsection{なぜ JRA-3Q か}

再解析にはいくつか選択肢がある。図\ref{fig:f12} に並べた。

\begin{figure}[H]
\centering
\includegraphics[width=1.0\linewidth]{compare/F12_reanalyses.pdf}
\caption[再解析プロダクトの比較]{主な再解析の比較。
左は横軸が水平格子間隔（\textbf{細かいほうが右}）、縦軸が気圧面プロダクトの上端高度。
本用途に適するのは右上で、下の帯は破裂高度 \SI{30}{\kilo\metre} に届かない範囲である。
右は対象期間と、気圧面の数・解析間隔。
\textbf{JRA-3Q の値だけが実測}で、本パイプラインが取得したファイルから読んでいる。
他は各提供元が公表している\textbf{気圧面プロダクト}の仕様を手で入力したもので、
モデル格子やモデル面の仕様ではない。
仕様は更新されうるので、自分の研究に選ぶときは最新の値を確認してほしい。}
\label{fig:f12}
\end{figure}

図から読めることが 3 つある。

\textbf{NCEP/NCAR R1 は使えない。} 気圧面プロダクトが \SI{10}{\hecto\pascal}
（約 \SI{31}{\kilo\metre}）で止まる。ちょうど破裂高度なので、
飛行のいちばん上の部分に風のデータが無いことになる。
格子も \SI{200}{\kilo\metre} 級で、対象領域に数点しか入らない。

\textbf{JRA-55 は解像度が足りない。} \ang{1.25} は和歌山の領域に 4 点しか置けない。
「場所による違い」を論じる本報告の目的に合わない。
2023 年で更新も止まっている。

\textbf{残るのは ERA5、MERRA-2、JRA-3Q の 3 つ}で、
どれも高さと解像度の条件は満たす。ここから先は優劣ではなく選択である。

JRA-3Q を選んだ理由を挙げる。

\begin{enumerate}
  \item \textbf{気象庁自身の製品である。}
        日本周辺の観測、特に国内のラジオゾンデや日本の静止衛星のデータを厚く同化している。
        対象領域が日本である以上、この点で他の再解析より精度が期待できる。
  \item \textbf{高さが足りる。}
        気圧面プロダクトが 45 面あり、上端は \SI{0.01}{\hecto\pascal}（約 \SI{80}{\kilo\metre}）に達する。
        スペースバルーンに必要な \SI{40}{\kilo\metre} までは十分に余裕がある。
        Open-Meteo の \SI{30}{\hecto\pascal} 上限という制約がここで解ける。
  \item \textbf{期間が長い。}
        1947 年 9 月から現在まで。今回は使わないが、将来もっと長い統計が要るときに困らない。
  \item \textbf{入手できる。}
        NSF NCAR の GDEX から無償で、しかも匿名で取得できる。
\end{enumerate}

\begin{whybox}{ERA5 を選ばなかったのは、性能の問題ではない}
図\ref{fig:f12} のとおり ERA5 は JRA-3Q より水平解像度が細かく（\SI{20}{\kilo\metre} 対
\SI{31}{\kilo\metre}）、時間解像度も 1 時間ごとで 6 倍細かい。
気圧面の上端は \SI{1}{\hecto\pascal}（約 \SI{48}{\kilo\metre}）で、
これもスペースバルーンには十分である。
本用途の要件だけを見れば ERA5 のほうが有利な点が多い。

選ばなかった理由は 2 つある。

第一に、\textbf{ERA5 は同じプロジェクトの別のメンバーが並行して解析している}。
2 人で同じ製品を重複して扱うより、別々の再解析を分担したほうが、
あとで突き合わせたときに「製品による違い」が見える。
本報告の数字が JRA-3Q 固有のものなのか再解析一般に言えることなのかは、
本報告だけでは判定できない。それは 2 つを比べて初めて分かる。

第二に、日本周辺という対象に対しては\ref{sec:data} 節の理由 1（気象庁自身の製品で、
国内の観測を厚く同化している）が有利に働く可能性がある。
これは\textbf{検証していない期待であって、測った事実ではない}。
どちらがこの領域で実際に正確かは、2 つの解析が出そろってから比べるべきことである。
\end{whybox}

\subsection{使ったプロダクトと設定}

\begin{table}[H]
\centering
\caption{取得したデータの仕様}
\begin{tabular}{ll}
\toprule
項目 & 値 \\
\midrule
データセット & JRA-3Q \texttt{anl\_p}（等圧面解析値）\\
配布元 & NSF NCAR GDEX, データセット \texttt{d640000} \\
水平解像度 & \ang{0.375}（$960 \times 480$ のガウス格子、約 \SI{40}{\kilo\metre}）\\
鉛直 & 気圧面 38 面（\SIrange{1000}{3}{\hecto\pascal}）\\
時間 & 6 時間ごと（00, 06, 12, 18 UTC）\\
期間 & 2016--2025 年（10 年、14{,}610 時刻）\\
変数 & 東西風 $u$、南北風 $v$、ジオポテンシャル高度 $Z$ \\
領域 & \ang{30}--\ang{48}N / \ang{132}--\ang{150}E から各地域を切り出し \\
\bottomrule
\end{tabular}
\end{table}

\paragraph{なぜ 2016--2025 年なのか}
10 年という長さは 2 つの要求の折り合いである。
短すぎると 1 つのビンに入る標本が減り、統計が暴れる。
長すぎると、前述の観測網の変化や、気候そのものの長期変化が混ざってくる。
また現実的な問題として、10 年分で取得は約 2.5 時間、キャッシュは約 \SI{5}{\giga\byte} である。

\paragraph{なぜ \ang{0.375} なのか}
JRA-3Q には \ang{1.25} の軽いプロダクトもあり、
北海道だけなら \ang{1.25} でも格子点が 54 個取れて図になる。
しかし和歌山の領域（\ang{2.5} $\times$ \ang{3}）には \ang{1.25} 格子だと点が 4 個しか入らず、
「場所による違い」を論じることができない。
\ang{0.375} なら和歌山でも約 60 点、北海道で約 540 点になる。

なお、解像度を上げても\textbf{分かることが増えるとは限らない}点には注意が要る。
JRA-3Q のモデル格子は約 \SI{40}{\kilo\metre} であり、
\ang{0.375} のプロダクトはその内挿である。
山脈 1 本のスケールの現象は、どちらの解像度でも解けていない。

\paragraph{なぜ気圧面なのか}
大気のモデルは地形に沿った座標で計算されるが、配布は気圧面か等高度面が普通である。
気圧面を選んだのは、スペースバルーンの高度計が実質的に気圧計であること、
そして既存のシミュレータが Open-Meteo の気圧面を前提に作られていて
換算の仕方を揃えられることによる。
気圧面から実際の高度への換算は第 \ref{sec:method} 章で述べる。

\paragraph{なぜ \SI{3}{\hecto\pascal} で切ったのか}
\SI{3}{\hecto\pascal} は約 \SI{40}{\kilo\metre} である。
ゴム気球の破裂高度はふつう \SIrange{25}{33}{\kilo\metre} なので、
これで飛行の全区間に余裕を持って収まる。
それより上の 7 面（\SI{2}{\hecto\pascal} 以上）は中間圏に入り、
スペースバルーンには関係がないうえ、図の縦軸を無駄に伸ばして対流圏を見づらくする。

\subsection{ライセンスと引用}

JRA-3Q は \textbf{CC BY-NC-SA 4.0} で提供されている \citep{gdex_d640000}。
これは次の 3 つを意味する。

\begin{itemize}
  \item \textbf{非商用に限る。} 商用目的では使えない。
  \item \textbf{継承。} 本報告や出力データは二次的著作物であり、
        再配布するなら同じ CC BY-NC-SA 4.0 を付けなければならない。
        リポジトリ本体のライセンスとは別である。
  \item \textbf{表示。} 出典を明記しなければならない。
\end{itemize}

学術・技術文書での引用は JRA-3Q の総合報告 \citep{kosaka2024jra3q} に対して行う。
この継承条項があるため、本パイプラインの出力は
Web アプリのビルドには混ぜず、独立したディレクトリに置いている。
